\documentstyle[%


righttag%


,11pt%


,amssymb%


,russian%               remove in Eng.


,izvru%                 Set izven% in Eng.


]{amsart}





% NB: \text{} constructions





\newcommand{\im}{\operatorname{Im}}


\newcommand{\re}{\operatorname{Re}}


\newcommand{\conv}{\operatorname{conv}}


\newcommand{\ind}{\operatornamewithlimits{ind}}


\newcommand{\tts}[1]{}





\theoremstyle{definition}


\theorembodyfont{\rm}


\newtheorem{cornonum}{\hskip\parindent Следствие}


\renewcommand{\thecornonum}{}





%\hoffset=-15mm%                  For Eng.


\setcounter{page}{60}


\god{2001}


\nomer{8~(471)} %                        Remove in Eng.


%\nomer{Vol.\,45, No.\,8}%               Set in Eng.


%\str{\pageref{begin}--\pageref{end}}%  Set in Eng.


\udk{517.982.2}





\author{\it Ю.Ф.\,Коробейник}


% NB:   ^^ remove in Eng.


\title{О сходимости рядов в локально выпуклых пространствах}


\thanks{Работа выполнена при финансовой поддержке Российского фонда


фундаментальных исследований (проект 99--01--01018).}





\date{}


%\pagestyle{myheadings}%         Set in Eng.


\pagestyle{plain} %              Remove in Eng.


\begin{document}


%\thispagestyle{plain}%          Set in Eng.


\maketitle


\label{begin}





{\bf 1.} Пусть $H$ --- локально выпуклое пространство над полем скаляров


$\Phi$ (всюду далее $\Phi=\Bbb C$ или $\Phi=\Bbb R$) с определяющим топологию


$\mu$ набором преднорм $Q=\{q\}$. Пусть, далее, $B$ --- некоторое множество


индексов и $x(t):B\to H$ --- отображение $B$ в $H$. Положим $X_B=\{x(t):


t\in B\}$. Пусть $\{t_k\}_{k=1}^{\infty}$ --- некоторая


последовательность индексов из $B$ и


$\Lambda :=\{t_k\}_{k=1}^{\infty}$, $X_{\Lambda}=\{x(t_k):


k\geqslant 1\}$. Всюду далее предполагается, что $q(x(t))>0$ $\forall


t\in B$, $\forall q\in Q$. Положим 


$\tau_q(t):=\ln q(x(t))$ $\forall t\in B$, $\forall q\in Q$.





Основная цель данной работы состоит в определении условий (по возможности,


типа критериев) сходимости ряда


\begin{equation}


\sum\limits_{k=1}^{\infty}c_kx(t_k),\quad c_k\in\Phi,\quad k=1,2,\dotsc


\end{equation}


Поскольку для любого сходящегося в $H$


ряда (1) выполняется необходимое условие сходимости


$\lim\limits_{k\to\infty}|c_k|


q(x(t_k))=0$ $\forall q\in Q$, то справедливо





\begin{prop}


Если ряд $(1)$ сходится в $H$, то


\begin{equation}


\lim\limits_{k\to\infty}\bigl(\ln|c_k|+\tau_q(t_k)\bigr)=-\infty\quad


\forall q\in Q.


\end{equation}


\end{prop}





\begin{cornonum}


Если ряд (1) сходится в $H$, то


\begin{equation}


\limsup\limits_{k\to\infty}(\ln|c_k|+\tau_q(t_k))<+\infty\ \; \forall q\in Q.


\end{equation}


\end{cornonum}





Будем говорить, что последовательность $\Lambda$ {\it ядерна} относительно


топологии $\mu$, если


\begin{equation}


\forall q\in Q\;\;\exists q_1\in Q:\;


\sum\limits_{k=1}^{\infty}


\exp\big(\tau_q(t_k)-\tau_{q_1}(t_k)\big)<\infty.


\end{equation}





\begin{prop}


Если $\Lambda$ ядерна относительно $\mu$ и


выполняется условие $(3)$, то ряд $(1)$ сходится абсолютно в $H$.


\end{prop}





Для доказательства достаточно зафиксировать произвольно $q\in Q$


и выбрать $q_1\in Q$ так, чтобы выполнялось условие (4). Тогда


в силу (3) и (4)


$$


\sum\limits_{k=1}^{\infty}|c_k|q(x(t_k))=


\sum\limits_{k=1}^{\infty}\exp\big(\ln|c_k|+\tau_{q_1}(t_k)\big)


\exp\big(\tau_q(t_k)-\tau_{q_1}(t_k)\big)<\infty.


$$





Из следствия предложения 1 и предложения 2 получаем





\begin{prop}


Для того чтобы ряд $(1)$ сходился $($или абсолютно


сходился$)$ в $H$, необходимо, а в случае, если последовательность $\Lambda$


ядерна относительно топологии $\mu$, и достаточно, чтобы выполнялось


условие $(3)$.


\end{prop}





Учитывая также предложение 1, получаем первый критерий сходимости ряда (1).





\begin{thm}


Пусть $\Lambda$ ядерна относительно $\mu$. Тогда


следующие утверждения равносильны\/$:$


\begin{itemize}


\item[1)] ряд $(1)$ сходится в $H;$


\item[2)] ряд $(1)$ абсолютно сходится в $H;$


\item[3)] выполняется условие $(2);$


\item[4)] выполняется условие $(3)$.


\end{itemize}


\end{thm}





Для проверки условия (4) ядерности нужно знать асимптотику поведения при


$k\to\infty$ величины $\tau_q(t_k)-\tau_{q_1}(t_k)$, что во многих случаях


довольно затруднительно. Однако нередко удается получить односторонние оценки


(сверху) указанной величины и, используя их, получить достаточные условия


ядерности $\Lambda$ относительно $\mu$. Дадим предварительно два определения.





Будем говорить, что топология $\mu$\; $\wt{d}$-{\it разделена},


если 


\begin{equation}


\forall q\in Q\ \; \exists q_1\in Q,\ \; \exists b<\infty,\ \;


\exists d_{q}(t): B\longrightarrow[0,+\infty)\mid


\tau_{q_1}(t) \geqslant \tau_q(t)+d_q(t)-b\ \, \forall t\in B.


\end{equation}


Далее, последовательность $\Lambda$\; $\wt{d}$-{\it ядерна} в $H$,


если топология $\mu$\; $\wt{d}$-разделена и функцию $d_q(t)$ в


условии (5) можно выбрать так, что


$$


\sum\limits_{k=1}^{\infty}


\exp\big(-d_{q}(t_k)\big)<\infty\ \; \forall q\in Q.


$$





Очевидным является





\begin{prop}


Если $\Lambda$\; $\wt{d}$-ядерна в $H$,


то $\Lambda$ ядерна относительно $\mu$.


\end{prop}





\begin{cornonum}


Если последовательность $\Lambda$\;


$\wt{d}$-ядерна в $H$, то утверждения 1)--4) теоремы 1 равносильны.


\end{cornonum}





Укажем теперь достаточно простые условия, обеспечивающие


$\wt{d}$-разделенность $\mu$ и $\wt{d}$-ядер\-ность


$\Lambda$ в $H$.


Hазовем последовательность


$\Lambda$ {\it сильно $\wt{d}$-ядерной} в $H$, если топология


$\mu$\; $\wt{d}$-разделена и если функцию $d_q(t)$ в (5) можно


выбрать так, чтобы


$$


\limsup\limits_{k\to\infty}\frac{\ln k}{d_q(t_k)}


<1\ \; \forall q\in Q.


$$


Eсли $\Lambda$ сильно $\wt{d}$--ядерна в $H$, то $\Lambda$ подавно


$\wt{d}$-ядерна в $H$.


Далее, будем говорить, что топология $\mu$\; $d$-разделена, если


$\mu$\; $\wt{d}$-разделена и функцию $d_q(t)$ в (5) можно


выбрать так, чтобы $d_q(t)=d(t)$ $\forall q\in Q$, $\forall t\in B$, где


$d(t):\;B\to[0,+\infty)$.


Пусть $\mu$\; $d$-разделена и пусть $q_1\in Q$, $d_q(t)=d(t)$ и постоянная


$b$  выбраны по $q\in Q$ так, чтобы выполнялось соотношение (5). Найдем теперь


по $q_1\in Q$ преднорму $q_2$ и постоянную $b_1$ так, чтобы 


$\tau_{q_2}(t)\geqslant \tau_{q_1}(t)+d_q(t)-b_1$


$\forall t\in B$. Тогда 


$\forall q\in Q$


$\exists q_2\in Q$, $\exists b_2\in\Bbb R$ $\forall t\in B$\;


$\tau_{q_2}(t)\geqslant \tau_{q}(t)+2d_q(t)-b_2$,


и вообще


$\forall q\in Q$,


$\forall N\geqslant 1$ $\exists q_N\in Q$, $\exists b_N\in\Bbb R$


$\forall t\in B$\;


$\tau_{q_N}(t)\geqslant \tau_{q}(t)+Nd_q(t)-b_N$.


Будем говорить, что $\Lambda$ имеет {\it конечный логарифмический тип}, если


$$


D_\Lambda:=\limsup\limits_{k\to\infty}\frac{\ln k}{d(t_k)}<\infty,


$$


и {\it нулевой логарифмический тип}, если $D_\Lambda=0$.





Из только что приведенных оценок сверху для $\tau_q(t)$ следует, что если


топология $\mu$\; $d$-разделена и $\Lambda$ имеет конечный логарифмический


$d$-тип, то $\Lambda$\; $\wt{d}$-ядерна в $H$.





Далее, скажем, что топология $\mu$ {\it слабо $d$-разделена}, если


$\mu$\; $\wt{d}$-разделена, причем функцию $d_q(t)$ в (5) можно


выбрать так, чтобы $d_q(t)=\alpha_qd(t)$ $\forall q\in Q$, $\forall t\in B$,


где $\alpha_q\in(0,+\infty)$. Легко показать, что если


$\mu$ слабо $d$-разделена и $\Lambda$ имеет нулевой логарифмический


$d$-тип, то $\Lambda$ сильно $\wt{d}$-ядерна в $H$.





Из приведенных соображений вытекает





\begin{thm}


Если $\Lambda$ сильно $\wt{d}$-ядерна в $H$, то


утверждения $1)$--\,$4)$ теоремы $1$ равносильны.


\end{thm}





\begin{cornonum}


Утверждения 1)--4) теоремы 1 равносильны, если выполнено


одно из двух следующих предположений:


\begin{itemize}


\item[(i)] $\mu$\; $d$-разделена, $D_\Lambda<+\infty$;


\item[(ii)] $\mu$ слабо $d$-разделена, $D_\Lambda=0$.


\end{itemize}


\end{cornonum}





{\bf 2.} Выведем теперь условия сходимости ряда (1) несколько иного вида.


Назовем последовательность $A=\bigl\{a_k\bigr\}_{k=1}^{\infty}$


положительных чисел {\it стандартной}, если $\lim\limits_{k\to\infty}a_k=


+\infty$. Из следствия предложения 1 получаем





\begin{prop}


Если $A$ --- любая стандартная последовательность,


и ряд $(1)$ сходится в $H$, то


\begin{equation}


 \limsup\limits_{k\to\infty}\frac{1}{a_k}


\big(\ln|c_k|+\tau_q(t_k)\big)\leqslant 0 \ \; \forall q\in Q.


\end{equation}


\end{prop}





Покажем, что при некоторых дополнительных предположениях условие типа (6)


обеспечивает абсолютную сходимость ряда (1).


Пусть $\Lambda$ ядерна относительно $\mu$. Будем говорить, что семейство


стандартных последовательностей $A(q)=\big\{a_k(q)\big\}_{k=1}^{\infty}$,


$q\in Q$, не нарушает ядерности $(\Lambda,\mu)$, если


\begin{equation}


\forall q\in Q\ \exists q_1\in Q,\ \exists \varepsilon_q>0:\ \;


\sum\limits_{k=1}^{\infty}\exp\big(\varepsilon_qa_k(q_1)+\tau_q(t_k)-


\tau_{q_1}(t_k)\big)<\infty.


\end{equation}





\begin{prop}


Пусть $\Lambda$ ядерна относительно $\mu$ и пусть семейство


стандартных последовательностей $A(q)$ не нарушает ядерности


$(\Lambda,\mu)$. Пусть, далее, выполнено условие


\begin{equation}


 \limsup\limits_{k\to\infty}\frac{1}{a_k(q)}


\bigl(\ln|c_k|+\tau_q(t_k)\bigr)\leqslant 0\ \; \forall q\in Q .


\end{equation}


Тогда ряд $(1)$ сходится абсолютно в $H$.


\end{prop}





Доказательство проводится так же, как и для предложения 2, но


с использованием (8) и (7).





\begin{cornonum}


Пусть $\Lambda$ ядерна относительно $\mu$ и пусть семейство


$A(q)$ не нарушает ядерности


$(\Lambda,\mu)$. Тогда любое из утверждений 1)--4) теоремы 1 равносильно


неравенству (8).


\end{cornonum}





Выясним, при каких условиях семейство стандартных последовательностей


$A(q)$ не нарушает ядерности $(\Lambda,\mu)$. Условие (7) равносильно тому,


что ряд Дирихле


$$


\sum\limits_{k=1}^{\infty}\beta_ke^{-a_k(q_1)x},\quad \beta_k=


\exp\bigl(\tau_q(t_k)-\tau_{q_1}(t_k)\bigr),\ \;k=1,2,\dotsc,


$$


имеет отрицательную абсциссу сходимости при любом $q\in Q$ и некотором


$q_1\in Q$. Пусть $A(q)$ таково, что


\begin{equation}


\lim_{k\to\infty}\frac{\ln k}{a_k(q)}=0\ \; \forall q\in Q .


\end{equation}


Тогда условие (7) равносильно тому, что


\begin{equation}


\forall q\in Q\ \; \exists q_1\in Q:\ \;


\limsup_{k\to\infty}


\frac{\tau_q(t_k)-\tau_{q_1}(t_k)}{a_k(q_1)}<0.


\end{equation}


Таким образом, если выполнены условия (9) и (10), то $A(q)$ не нарушает


ядерности $(\Lambda,\mu)$.





Точно так же показываем при условии (9), что если $\mu$\;


$\wt{d}$-разделена и


$$


\forall q\in Q\ \; \exists q_1\in Q: \ \;


\limsup_{k\to\infty}


\frac{-d_q(t_k)}{a_k(q_1)}<0,


$$


т.\,е. если


\begin{equation}


\forall q\in Q\ \; \exists q_1\in Q:\ \;


\limsup_{k\to\infty} \frac{a_k(q_1)}{d_q(t_k)}<\infty,


\end{equation}


то $A(q)$ не нарушает ядерности $(\Lambda,\mu)$.


Следовательно, справедлива





\begin{thm}


Пусть $\Lambda$\; $\wt{d}$-ядерна в $H$ и семейство


стандартных последовательностей $A(q)$ удовлетворяет условиям $(9)$, $(11)$


или $(9)$, $(10)$. Тогда любое из утверждений $1)$--\,$4)$


теоремы $1$ равносильно неравенству $(8)$.


\end{thm}





Рассмотрим некоторые более специальные ситуации, считая, что условие (9)


выполнено. Пусть сначала справедливо предположение (i). Пусть еще


последовательность $A=\bigl\{a_k\bigr\}_{k=1}^{\infty}$ такова, что


\begin{equation}


\limsup\limits_{k\to\infty}\dfrac{a_k}{d(t_k)}<\infty,


\end{equation}


а коэффициенты $c_k$ ряда(1) удовлетворяют условию (6). Покажем, что тогда


ряд (1) абсолютно сходится. Согласно сделанным предположениям


$$


\exists T_1<\infty,\;\ \exists T_2<\infty\;\ \forall k\geqslant 1\ \;


a_k\leqslant T_1d(t_k),\quad \ln(1+k)\leqslant T_2d(t_k).


$$


Зафиксируем произвольно $q\in Q$ и выберем число $N>T_2+1$. Найдем $q_1\in Q$


и $b<\infty$ так, чтобы $\tau_{q_1}(t)\geqslant \tau_q(t)


+Nd(t)-b$ $\forall t\in B$.


Выберем $k_0\geqslant 1$ настолько большим, чтобы 


$\ln|c_k|+\tau_{q_1}(t_k)<a_k/T_1$ $\forall k \geqslant k_0$.


Тогда 


$$


|c_k|\exp \tau_q(t_k)\leqslant \exp\bigl(d(t_k)-Nd(t_k)+b\bigr)\leqslant


b_1(1+k)^{-\frac{N-1}{T_2}}\ \; \forall k\geqslant k_0.


$$





Таким же образом доказываются следующие утверждения.





\begin{prop}


Пусть имеют место предположения $(\mathrm i)$ и, кроме того,


\begin{gather}


\lim\limits_{k\to\infty}\dfrac{a_k}{d(t_k)}=0;\\


%


 \limsup\limits_{k\to \infty}\dfrac{1}{a_k}\bigl(


\ln|c_k|+\tau_q(t_k)\bigr)<\infty\ \; \forall q\in Q.


\end{gather}


Тогда ряд $(1)$ сходится абсолютно в $H$.


\end{prop}





\begin{prop}


Пусть для стандартной последовательности $A$


выполняется условие $(12)$, и, кроме того, имеет место предположение


$(\mathrm{ii})$. Тогда, если коэффициенты ряда $(1)$ удовлетворяют


условию $(6)$, то ряд сходится абсолютно в $H$.


\end{prop}





\begin{prop}


Пусть $A$ удовлетворяет условию $(13)$, и пусть выполнены предположения


$(\mathrm{ii})$. Тогда, если коэффициенты ряда $(1)$


удовлетворяют условию $(14)$, то ряд сходится


абсолютно в $H$.


\end{prop}





Из полученных результатов следуют две теоремы.





\begin{thm}


Пусть стандартная последовательность $A$ удовлетворяет


условию $(12)$, и пусть выполнено хотя бы одно из предположений


$(\mathrm i)$, $(\mathrm{ii})$.


Тогда любое из утверждений\linebreak $1)$--\,$4)$ теоремы $1$


равносильно соотношению $(6)$.


\end{thm}





\begin{thm}


Пусть стандартная последовательность $A$ удовлетворяет


условию $(13)$, и пусть выполнено хотя бы одно из предположений


$(\mathrm i)$, $(\mathrm{ii})$.


Тогда любое из утверждений\linebreak


$1)$--\,$4)$ теоремы $1$ равносильно каждому из


неравенств $(6)$, $(14)$.


\end{thm}





{\bf 3.} Все приведенные выше результаты применимы, в частности, к


пространству Фреше, топология $\mu$ в котором, как известно, задается не более


чем счетным набором преднорм $Q=\{q_n:n=1,2,\dotsc\}$. Читатель без труда


получит формулировки вышеприведенных общих результатов в этом важном частном


случае.





К сожалению, во многих локально выпуклых пространствах преднормы, определяющие


топологию, не допускают простого описания и, как правило, топология в них


описывается иным способом. Это относится, в частности, к играющим важную роль


в анализе $IF$-пространствам, т.\,е. отделимым внутренним индуктивным пределам


неубывающей последовательности пространств Фреше. Однако к этому классу


пространств удается при одном дополнительном предположении применить


полученные выше результаты. Следуя [1], скажем, что $IF$-пространство


$H=\ind H_n$ обладает свойством ($Y$), если для любой последовательности


$(v_n)_{n=1}^{\infty}$ элементов из $H$, сходящейся в $H$ к


некоторому элементу $v\in H$, найдется номер $m$ такой, что $v\in H_m$,


$v_k\in H_m$ $\forall k\geqslant 1$ и $v_k\to v$ в $H_m$. Условимся


называть любое такое $IF$-пространство  $(IF)_0$-пространством.





Как известно [1], к $(IF)_0$-пространствам относятся $LN^*$-пространства,


т.\,е. $IF$-пространства, в которых каждое $H_n$ вложено вполне непрерывно в


$H_{n+1}$, а также $LF$-пространства, т.\,е. строгие индуктивные пределы


последовательности пространств Фреше.





Всюду далее в этом параграфе $H=\ind H_n$ --- $(IF)_0$-пространство;


$H_n$ --- пространство Фреше с топологией $\mu_n$, заданной набором преднорм


$Q_n=(q_{n,k})_{k=1}^{\infty}$; $H_n\hookrightarrow H_{n+1}$,


$n=1,2,\dotsc$; $x(t)\in H_1$. Положим 


$\ln q_{m,n}(x(t))=:\tau_{m,n}(t)$ $\forall t\in B$,


$m,n=1,2, \dotsc$





\begin{thm}


Если последовательность


$\Lambda=(t_k)_{k=1}^{\infty}$ ядерна относительно каждой


топологии $\mu_n$, то следующие утверждения равносильны\/$:$


\begin{itemize}


\item[1)] ряд $(1)$ сходится в $H;$


\item[2)] $\exists p\geqslant 1:$ ряд $(1)$ сходится в $H_p;$


\item[3)] ряд $(1)$ сходится абсолютно в $H;$


\item[4)] $\exists m\geqslant 1:$  ряд $(1)$ сходится абсолютно в $H_m;$


\item[5)] $\exists j\geqslant 1:$


$\forall m\geqslant 1$ $\lim\limits_{k\to\infty}


\big(\ln|c_k|+\tau_{j,m}(t_k)\big)=-\infty;$


\item[6)] $\exists r\geqslant 1:$


$\forall n\geqslant 1:$ $\limsup\limits_{k\to\infty}


\bigl(\ln|c_k|+\tau_{r,n}(t_k)\bigr)<\infty$.


\end{itemize}


\end{thm}





\begin{pf}


Так как $H$ обладает свойством ($Y$), то


$1)\Longleftrightarrow 2)$. Далее, всегда $4)\Longrightarrow 3)


\Longrightarrow 1)$; $2)\Longrightarrow 5) \Longrightarrow 6)$.


Наконец, $6)\Longrightarrow 4)$  по


предложению 2 (с $m=r$).


\end{pf}





Учитывая также следствие предложения 4 и теорему 2, из теоремы 6 получаем





\begin{cor}


Если последовательность


$\Lambda$\; $\wt d_n$-ядерна в $H$ при любом $n\geqslant 1$, то


утверждения 1)--6) теоремы 6 равносильны.


\end{cor}





\begin{cor}


Пусть выполнено хотя бы одно из двух предположений:


\begin{itemize}


\item[(i)$_0$] $\forall n\geqslant 1$ топология $\mu_n$\;


$d_n$-разделена и $\Lambda$


имеет конечный логарифмический $d_n$-тип;


\item[(ii)$_0$]


$\forall n\geqslant 1$ топология $\mu_n$\; слабо $d_n$-разделена и


$\Lambda$ имеет нулевой логарифмический $d_n$-тип.


\end{itemize}


Тогда утверждения 1)--6) теоремы 6 равносильны.


\end{cor}





Переходим теперь к критериям сходимости, в которых участвует стандартная


стационарная


последовательность $A=\{a_k\}_{k=1}^{\infty}$. Положим $d_{q_{k,n}}


(t):=d_{k,n}(t)$ $\forall k,n\geqslant 1$, $\forall t\in B$. Примерно так


же, как теорема 6, но с помощью теоремы 3 доказывается





\begin{thm}


Пусть $\Lambda$\; $\wt d_n$-ядерна в $H_n$ $\forall n\geqslant 1$.


Пусть, далее, $A$ удовлетворяет условию $(9)$, а также


соотношению


\begin{equation}


\limsup\limits_{k\to\infty}\dfrac{a_k}


{d_{j,n}(t_k)}<\infty\ \; \forall j,n\geqslant 1.


\end{equation}


Тогда любое из утверждений $1)$--\,$6)$ теоремы $6$ равносильно неравенствам


\begin{equation}


\exists n\geqslant 1:\ \;\forall j\geqslant 1 \quad


\limsup\limits_{k\to\infty}\dfrac{1}{a_k}


\bigl(\ln|c_k|+\tau_{n,j}(t_k)\bigr)\leqslant 0.


\end{equation}


\end{thm}





Далее, из теорем 4 и 5 следуют две теоремы.





\begin{thm}


Пусть стандартная последовательность $A$ удовлетворяет


условию $(15)$ и хотя бы одному из предположений $(\mathrm i)_0$,


$(\mathrm{ii})_0$. Тогда


любое из утверждений $1)$--\,$6)$ теоремы $6$ равносильно неравенствам $(16)$.


\end{thm}





\begin{thm}


Пусть выполняется хотя бы одно из


предположений $(\mathrm i)_0$, $(\mathrm{ii})_0$,


и пусть стандартная стационарная


последовательность $A$


такова, что $\lim\limits_{k\to\infty}a_k/d_n(t_k)=0$ $\forall n\geqslant 1$.


Тогда соотношение


$$


\exists r\geqslant 1:\ \; \forall p\geqslant 1 \ \;


\limsup\limits_{k\to\infty}\dfrac{1}{a_k}


\bigl(\ln|c_k|+\tau_{r,p}(t_k)\bigr)<\infty


$$


равносильно неравенствам $(16)$, а также


любому из утверждений $1)$--\,$6)$ теоремы $6$.


\end{thm}





{\bf 4.} Рассмотрим более специальную ситуацию, когда


$(IF)_0$-пространство является $H=\ind H_n$,


где $\forall n\geqslant 1\;$ $H_n$ --- $B$-пространство с нормой


$\|\cdot\|_n$, вполне непрерывно вложенное в $B_{n+1}$.





Если ряд (1) сходится в $H$, то $\exists n\geqslant 1$: ряд (1) сходится в


$H_n$. Отсюда


\begin{equation}


\exists n\geqslant 1:\lim\limits_{k\to\infty}


\bigl(\ln|c_k|+\ln\|x(t_k)\|_n\bigr)=-\infty,


\end{equation}


и


\begin{equation}


\exists n\geqslant 1:\limsup\limits_{k\to\infty}


\bigl(\ln|c_k|+\ln\|x(t_k)\|_n\bigr)<\infty.


\end{equation}





Будем считать, что 


$\|x(t_k)\|>0$ $\forall n\geqslant 1$, $\forall k\geqslant 1$,


и назовем последовательность $\Lambda$


{\it индуктивно-ядерной} в $H$, если


$$


\forall n\geqslant 1\ \; \exists m\geqslant 1:\ \;


\sum\limits_{k=1}^{\infty}\dfrac{\|x(t_k)\|_m}{\|x(t_k)\|_n}<\infty.


$$


Тогда, если выполнено условие (18) и $\Lambda$ индуктивно ядерна в $H$,


то


$$


\sum\limits_{k=1}^{\infty}|c_k|\,\|x(t_k)\|_m=


\sum\limits_{k=1}^{\infty}|c_k|\,\|x(t_k)\|_n\,


\dfrac{\|x(t_k)\|_m}{\|x(t_k)\|_n}<\infty,


$$


и ряд (1) абсолютно сходится в $H$. Таким образом, справедливо





\begin{prop}


Если ряд $(1)$ сходится в $H$, то выполнено условие


$(17)$. Обратно, если $\Lambda$ индуктивно ядерна в $H$ и выполнено условие


$(18)$, то ряд $(1)$ сходится абсолютно в $H$.


\end{prop}





{\bf 5.}


Рассмотрим некоторые примеры, уделив основное внимание рядам экспонент. Пусть


$B\subseteq \Bbb C^p$, $t=(t_1,\dotsc,t_p)\in B$, $p\geqslant 1$,


$z=(z_1,\dotsc,z_p)\in\Bbb C^p$ или $z\in\Bbb R^p$, $\langle z,t\rangle=


\sum\limits_{k=1}^{p}z_kt_k$; $|t|_p:=|t_1|+\dotsb+|t_p|$. Пусть сначала


$0<\sigma<\infty$, $1<\rho<\infty$, $y(z)$ --- любая целая функция в


$\Bbb C^p$ и $\|y\|^\sigma_\rho:=\sup\limits_{r>0}M_r(y)/


\exp(\sigma r^\rho)$, где $M_r(y)=\max\{|y(z)|:\;|z_k|\leqslant


r,\ k=1,2,\dotsc, p\}$. Положим $\forall t\in\Bbb C^p$ $x(t)=\exp\langle


z,t\rangle$. Тогда, как известно,


$$


\ln\|x(t)\|_\rho^\sigma=(|t|_p)^{\Tilde\rho}


A_\rho^{(\sigma)}, \quad


\text{где} \quad \wt\rho=\dfrac{\rho}{\rho-1};\quad


\dfrac{1}{A_\rho^{(\sigma)}}=\wt\rho(\rho\sigma)^{\Tilde\rho/\rho}.


$$





1. Пусть $H=[\rho,\infty]_p$ --- пространство Фреше всех целых функций в


$\Bbb C^p$ порядка не выше, чем $\rho$, с определяющим топологию набором


норм


$$


\|y\|_{n,\rho}:=\|y\|_{\rho_n^0}^{1}, \quad \text{где} \quad


\rho_n^0=\rho\bigg(1+\dfrac 1n\bigg),\quad n=1,2,\dotsc


$$


В этом случае $\forall t\in\Bbb C$


$$


\ln\|x(t)\|_{n,\rho}=A^{(1)}_{\rho_n^0}


(|t|_p)^{\rho_n},\quad \text{где} \quad \rho_n=


\dfrac{\rho_n^0}{\rho_n^0-1}.


$$


Заметим, что $\rho_n\uparrow\rho/(\rho-1)$ при $\rho>1$ и $\rho_n


\uparrow +\infty$ при $\rho=1$. Будем далее считать, что


$\lim\limits_{k\to\infty}|t_{(k)}|_p=\infty$,


где $t_{(k)}=(t_{k,1},t_{k,2},\dotsc, t_{k,p})$. Более того, будем


предполагать, что числа $|t_{(k)}|_p$ возрастают не слишком медленно,


а именно,


\begin{equation}


\limsup\limits_{k\to\infty}\dfrac{\ln\ln k}{\ln |t_{(k)}|_p}<\infty.


\end{equation}


Из условия (19) следует, что $\exists m_0\geqslant 1$:


\begin{equation}


\lim\limits_{k\to\infty}\dfrac{\ln k}{|t_{(k)}|_p^{m_0}}=0.


\end{equation}


Согласно общей теории рядов Дирихле (напр., [2], гл.\,IV, \S\,1)


ряд $\sum\limits_{k=1}^{\infty}\exp\bigl(-x\lambda_k)$, где $\lambda_k=


|t_{(k)}|_p^{m_0}$, при условии (20) сходится при всех $x>0$ .





Далее, 


$$


\forall n\geqslant m_0,\; \ \forall D<\infty, \forall m>n\ \;


\exists k_0\geqslant 1\ \; \forall k\geqslant k_0\quad


D\tau_n(t_{(k)})-\tau_m(t_{(k)})<-\tau_n(t_{(k)}).


$$


Следовательно, условие (4) выполняется, и $\Lambda=(t_k)_{k=1}^{\infty}$


ядерна относительно топологии $[\rho,\infty]_p$. Кроме того, если


в качестве семейства стандартных последовательностей $A(n)=\{


a_{n,k}\}_{k=1}^{\infty}$ ($n=1,2,\dotsc$) взять 


$A(n)=\{(|t_{(k)}|_p)^{\rho_n}\}_{k=1}^{\infty}$ $\forall n\geqslant 1$, то


условие (7) также выполняется, и $A(n)$ не нарушает ядерности $(\Lambda,


\mu)$. Непосредственно из теоремы 1 и следствия предложения 6 получаем, что


для ряда


\begin{equation}


\sum\limits_{k=1}^{\infty}c_k\exp\langle t_{(k)},z\rangle


\end{equation}


равносильны утверждения:


\begin{itemize}


\item[1)] ряд (21) сходится в $[\rho,\infty]_p$;


\item[2)] ряд (21) сходится абсолютно в $[\rho,\infty]_p$;


\item[3)] $\limsup\limits_{k\to\infty}\dfrac{\ln|c_k|}{|t_{(k)}|_p^{\rho_n}}


\leqslant -A^{(1)}_{\rho_n^0}, \ \; n\geqslant m_0$.


\end{itemize}


Нетрудно проверить, что условие 3) равносильно тому, что


$$


\lim\limits_{k\to\infty}\dfrac{\ln|c_k|}{|t_{(k)}|_p^{\rho_n}}


= -\infty \ \; \forall n\geqslant m_0 .


$$


В свою очередь, последнее условие равносильно при $\rho=1$ соотношениям


$$


\lim\limits_{k\to\infty}\dfrac{\ln|c_k|}{|t_{(k)}|_p^N}


= -\infty\ \; \forall N<\infty,


$$


а при $\rho>1$ равносильно равенствам


$$


\lim\limits_{k\to\infty}\dfrac{\ln|c_k|}


{|t_{(k)}|_p^{\Tilde\rho-\varepsilon}}= -\infty\ \; \forall \varepsilon>0.


$$





2. Пусть теперь


$H=[\rho,\sigma]_p$ --- пространство всех тех целых функций в


$\Bbb C^p$, у которых тип при порядке $\rho$ (по совокупности переменных)


не превышает $\sigma$, $1<\rho<\infty$, $0\leqslant \sigma<\infty$.


Топология Фреше в $[\rho,\sigma]_p$ задается счетным набором норм


$\|y\|^{\sigma_n}_{\rho}$, где $\sigma _n=\sigma+1/n$


$\forall n\geqslant 1$. В данном случае $\forall t\in \Bbb C^p$, как выше,


$$


\ln\|x(t)\|^{\sigma_n}_{\rho}=\bigl(|t|_p\bigr)^{\Tilde\rho}


A^{(\sigma+1/n)}_{\rho},\quad n=1,2,\dotsc


$$


При этом


$$


0<A^{(\sigma+1/n)}_{\rho}\uparrow\dfrac{1}{\wt\rho}


(\rho)^{-\wt\rho/\rho}(\sigma)^{-\wt\rho/\rho}=


A^{(\sigma)}_{\rho}.


$$


Имеем 


$$


\tau_n(t)-\tau_m(t)=(|t|_p)^{\Tilde\rho}


(A^{(\sigma+1/n)}_{\rho}-A^{(\sigma+1/m)}_{\rho})\leqslant


-(A^{(\sigma+1/(n+1))}_{\rho}-A^{(\sigma+1/n)}_{\rho})


|t|_p^{\Tilde\rho}


\ \forall n\geqslant 1,\  \forall m>n,\  \forall t\in\Bbb C^p.


$$





Если в качестве $d(t)$ взять $|t|_p^{\Tilde\rho}$, то топология


$[\rho,\sigma]_p$ слабо $d$-разделена. Далее, если


\begin{equation}


\lim\limits_{k\to\infty}\dfrac{\ln(1+k)}


{|t_{(k)}|_p^{\Tilde\rho}}= 0,


\end{equation}


то $\Lambda=\bigl(t_{(k)}\bigr)^{\infty}_{k=1}$ имеет нулевой логарифмический


тип. Согласно теореме 2 равносильны утверждения:


\begin{itemize}


\item[1)] ряд (21) сходится в $[\rho,\sigma]_p$;


\item[2)] ряд (21) сходится абсолютно в $[\rho,\sigma]_p$;


\item[3)] $\lim\limits_{k\to\infty}(\ln|c_k|+


A^{(\sigma+1/n)}_{\rho}(|t_{(k)}|_p)^{\Tilde\rho})=


-\infty$ $\forall n\geqslant 1$ ;


\item[4)] $\limsup\limits_{k\to\infty}(\ln|c_k|+


A^{(\sigma+1/n)}_{\rho}|t_{(k)}|_p^{\Tilde\rho})<\infty$


$\forall n\geqslant 1$.


\end{itemize}


Взяв в качестве стандартной последовательности


$\{(|t_{(k)}|_p)^{\Tilde\rho}\}^{\infty}_{k=1}$


и учитывая условие (22), получаем непосредственно из теоремы 4, что любое из


утверждений 1)--4) равносильно тому, что


\begin{itemize}


\item[5)] $\limsup\limits_{k\to\infty}\dfrac{\ln|c_k|}


{\bigl(|t_{(k)}|_p\bigr)^{\Tilde\rho}}\leqslant - A^{(\sigma+1/n)}_{\rho}$


$\forall n\geqslant 1$.


\end{itemize}


В свою очередь, утверждение 5) равносильно условию


\begin{equation}


\text{при}\ \sigma>0\ \limsup\limits_{k\to\infty}\dfrac{\ln|c_k|}


{(|t_{(k)}|_p)^{\Tilde\rho}}\leqslant - A^{\sigma}_{\rho};


\ \;


\text{при}\ \sigma=0\ 


\lim\limits_{k\to\infty}\dfrac{\ln|c_k|}


{(|t_{(k)}|_p)^{\rho}}= - \infty,


\end{equation}


которое является наиболее простым и легко проверяемым критерием


сходимости (или абсолютной сходимости) ряда (21) в $[\rho,\sigma]_p$


из четырех равносильных утверждений 3)--5), (23) при $\sigma\in[0,+\infty)$.


\medskip





3. Пусть $1<\rho<\infty$ и $H=[\rho,\infty)_n$ --- пространство всех целых


функций конечного типа при порядке $\rho$ (по совокупности переменных). Здесь


$[\rho,\infty)_n=\ind\limits_{m}H_m$, $H_m=[\rho,m]$, $m=1,2,\dotsc$, и


$H$ --- $LN^*$-пространство (следовательно, $(IF)_0$-пространство). Пусть


выполнено условие (22). Имеем 


$\tau_{m,n}=(|t|_p)^{\Tilde \rho}A_{\rho}^{m+1/n}$ $\forall m,n\geqslant 1$,


$t\in\Bbb C^p$,


где числа $A^\sigma_\rho$ определены выше. Отсюда $\forall


m\geqslant 1$ топология $\mu_m$ в $H_m$ слабо $d_m$-разделена, где


$d_m(t)=d(t)=\bigl(|t|_p\bigr)^{\Tilde \rho}$. Учитывая также условие


(22) и взяв в качестве стандартной последовательности $A=(a_k)


_{k=1}^{\infty}$, $a_k=(|t_{(k)}|_p)^{\Tilde \rho}$, согласно теореме~8


заключаем, что соотношение


\begin{equation}


\exists n\geqslant 1\ \; \forall j\geqslant 1\ \;


\limsup\limits_{k\to\infty}\dfrac{\ln|c_k|}{(|t_{(k)}|_p)^


{\wt\rho}}\leqslant -A_\rho^{n+1/j}


\end{equation}


равносильно любому из утверждений 1)--6) теоремы 6, в которых


$H=[\rho,\infty)_p$, $H_m=[\rho,m]_p$,


$\tau_{m,n}(t)=(|t|_p)^{\Tilde \rho}A_{\rho}^{m+1/n}$.





Нетрудно видеть, что условие (24) эквивалентно тому, что


$$


\exists n\geqslant 1:\ \;


\limsup\limits_{k\to\infty}\dfrac{\ln|c_k|}{(|t_{(k)}|_p)^


{\wt\rho}}\leqslant -A_\rho^n.


$$


В свою очередь, т.\,к. $\lim\limits_{n\to\infty}A^n_\rho=0$, то последнее


неравенство равносильно тому, что


\begin{equation}


\limsup\limits_{k\to\infty}\dfrac{\ln|c_k|}{(|t_{(k)}|_p)^


{\wt\rho}}<0.


\end{equation}


Очевидно, неравенство (25) и доставляет наиболее простой критерий


сходимости (а также абсолютной сходимости) ряда (21) в пространстве


$[\rho,\infty)_p$.


\medskip





4. Пусть $H=\cal H(\Bbb C^p)$ --- пространство Фреше всех целых функций в


$\Bbb C^p$ со счетным набором норм: $\|y\|_n=\max\{|y(v)|:|v_k|\leqslant


n,\ k=1,2,\dotsc,p\}$, $n=1,2,\dotsc$; $x(t)=\exp\langle z,t\rangle$,


$t\in\Bbb C^p$. В данном случае 


$\ln\|x(t)\|_n=|t|_p n$ $\forall n\geqslant 1$,


и топология $d$-разделена с $d(t)=|t|_p$.


Пусть $\Lambda=(|t_{(k)}|_p)$ и $\Lambda$ имеет конечный


логарифмический тип, т.\,е.


$$


\limsup\limits_{k\to\infty}\dfrac{\ln(1+k)}{|t_{(k)}|_p}<\infty.


$$


При $a_k=|t_{(k)}|_p$ согласно теореме 4 ряд (21) сходится (или


абсолютно сходится) в $\cal H(\Bbb C^p)$ тогда и только тогда, когда


$$


\limsup\limits_{k\to\infty}\dfrac{\ln|c_k|}{|t_{(k)}|_p}\leqslant -n


\ \;\forall n\geqslant 1,


$$


т.\,е. когда


$\limsup\limits_{k\to\infty}\dfrac{\ln|c_k|}{|t_{(k)}|_p}=-\infty.$


\medskip





5. Пусть $G$ --- содержащая начало координат ограниченная выпуклая область


в $\Bbb C^p$ с опорной функцией $h_G(t)=\sup\limits_{z\in G}\re\langle


z,t\rangle$, и пусть $H=H(G)$ --- пространство Фреше с набором норм


$$


\|y\|_n=\max\{|y(z)|:z\in\overline{q_nG}\},\quad n=1,2,\dotsc,\quad


0<q_n\uparrow 1.


$$


Если, как выше, $x(t)=\exp\langle z,t\rangle$, то $\ln\|x(t)\|_n=q_nh_G(t)$


$\forall n\geqslant 1$,


$\forall t\in\Bbb C^p$. В данном случае топология


в $H(G)$ слабо $d$-разделена с $d(t)=h_G(t)$. Пусть


$\Lambda=\bigl(t_{(k)}\bigr)_{k=1}^{\infty}$ и $\Lambda$ имеет нулевой


логарифмический тип, т.\,е.


\begin{equation}


\lim\limits_{k\to\infty}\dfrac{\ln(1+k)}{h_G(t_{(k)})}=0.


\end{equation}


По теореме 4 при $a_k=h_G(t_{(k)})$ ряд (21) сходится (или абсолютно


сходится) в $H(G)$ тогда и только тогда, когда


$$


\limsup\limits_{k\to\infty}\dfrac{\ln|c_k|}{h_G(t_{(k)})}\leqslant -q_n


\ \; \forall n\geqslant 1,


$$


или, что равносильно,


\begin{equation}


\limsup\limits_{k\to\infty}\dfrac{\ln|c_k|}{h_G(t_{(k)})}\leqslant -1.


\end{equation}


Заметим, что условие (26) равносильно тому, что


\begin{equation}


\lim\limits_{k\to\infty}\dfrac{\ln(1+k)}{|t_{(k)}|_p}=0.


\end{equation}





6. Аналогично, если $0\in G$ и $G$ --- выпуклая область (уже не обязательно


ограниченная), $H=H(G)$, $\|y\|_n=\max\{|y(z)|:z\in F_n\}$, $n=1,2,\dotsc$,


где $0\in F_n\subset F_{n+1}^0\subset G$ и $F_n$ --- компакт,


то $\ln\|x(t)\|_n=h_{F_n}(t)$ $\forall t\in\Bbb C^p$, $\forall n\geqslant 1$.


Здесь топология в $H(G)$ слабо $\wt d$-разделена, причем $\forall n\geqslant


1$\; $d_n(t)=h_{F_n}(t)$: для любого $n\geqslant 1$ найдутся достаточно


большое $m$ и малое $\varepsilon_n>0$ такие, что $\ln\|x(t)\|_n-


\ln\|x(t)\|_m\leqslant -\varepsilon_n\ln\|x(t)\|_n$ $\forall t\in\Bbb


C^p$. Полагая $a_{k,n}=h_{F_n}(t_k)$ $\forall n\geqslant 1$ , $k=1,2,


\dotsc$, получаем (учитывая неравенство $h_{F_{n+1}}(t)/h_{F_{n}}(t)


\geqslant b_n>1$\; $\forall t\in\Bbb C^p$)


согласно теореме 3, что ряд (21) при


условии (28) сходится (или абсолютно сходится) в $H(G)$ тогда и только


тогда, когда


\begin{equation}


\lim\limits_{k\to\infty}\dfrac{\ln|c_k|}{h_{F_n}(t_{(k)})}


\leqslant -1\ \;


\forall n\geqslant 1.


\end{equation}





7. Пусть $0\in F$, $F$ --- выпуклый компакт в $\Bbb C^p$ и $H=H(F)$ ---


пространство аналитических ростков на $F$ с индуктивной топологией


$\ind\limits_{n}H(G_n)$, где $F\subset G_n$ $\forall n\geqslant 1$, $G_n$ ---


выпуклая ограниченная область в $\Bbb C^p$, $\overline{G_{n+1}}\subset G_n$.


Так как $H(F)$ --- $(IF)_0$-пространство, то ряд (21) при условии (28)


сходится (или абсолютно сходится) в $H(F)$ тогда и только тогда, когда он


сходится (или абсолютно сходится) в некотором $H(G_n)$, т.\,е. когда (в силу


(27))


\begin{equation}


\exists n\geqslant 1\ \; \forall m\geqslant 1\ \;


\limsup\limits_{k\to\infty}\dfrac{\ln|c_k|}{h_{G_n}(t_{(k)})}


\leqslant -q_m.


\end{equation}


Здесь в качестве компактов $F_{n,m}$, аппроксимирующих изнутри область $G_n$,


берутся $q_m\overline{G_n}$, где $0<q_m\uparrow 1$. Очевидно, условие


(30) равносильно такому критерию сходимости


$$


\exists n\geqslant 1:\ \;


\limsup\limits_{k\to\infty}\dfrac{\ln|c_k|}{h_{G_n}(t_{(k)})}


\leqslant -1.


$$


Если, в частности, $F=\{0\}$, то можно положить $G_n=r_n\cal E_n$, где


$r_n\downarrow 0$, $\cal E_1=\{z:|z_k|<1,\ k=1,2,\dotsc,p\}$. Тогда $h_{G_n}


(t_{(k)})=r_n|t_{(k)}|$, и ряд (21) при условии (28) сходится (абсолютно


сходится) в $H(\{0\})$ тогда и только тогда, когда


$$


\limsup\limits_{k\to\infty}\dfrac{\ln|c_k|}{|t_{(k)}|_p}<0.


$$





8. Пусть $F$ --- выпуклый компакт в $\Bbb R^p$ с непустой внутренностью


$F^0$, $0\in F$, и пусть $C^\infty(F)$ --- пространство Фреше всех


комплекснозначных функций $y(x):F\to\Bbb C$, имеющих в $F^0$ частные


производные всех порядков по всем $p$ переменным $x_1,\dots,x_p$,


равномерно непрерывные


в $F^0$. Топология в $C^\infty (F)$ задается набором норм


$$


\|y\|_n=\sup\bigg\{\bigg|\dfrac{\partial^{|\alpha|_p}y(x)}


{\partial x_1^{\alpha_1}\dotsb \partial x_p^{\alpha_p}}\bigg|:|\alpha|_p


\leqslant n,\ x\in F_0\bigg\}, \quad n=0,1,2,\dotsc


$$


(здесь


$\alpha=(\alpha_1,\dots,\alpha_p)$). Имеем 


$$


\ln\|x(t)\|_n=h_F(t)+\ln\max\{|t_1|^{\alpha_1}\dotsb |t_p|^{\alpha_p}:


|\alpha|_p\leqslant n\}=\gamma_n(t)+h_F(t)


\ \; \forall t\in\Bbb C^p\ \; \text{и}\ \;


x(t)=\exp\langle t,x\rangle.


$$


Легко проверить, что


\begin{gather*}


\max\{|t_1|^{\alpha_1}\dotsb |t_p|^{\alpha_p}:|\alpha|_p\leqslant n+1\}


\geqslant\\


\geqslant


\max\{|t_1|^{\alpha_1}\dotsb |t_p|^{\alpha_p}:|\alpha|_p\leqslant n\}


\max\{|t_1|,|t_2|,\dotsc,|t_p|\} \geqslant \\


\geqslant


\max\{|t_1|^{\alpha_1}\dotsb |t_p|^{\alpha_p}:|\alpha|_p\leqslant n\}


\dfrac{|t|_p}{p}.


\end{gather*}


Отсюда $\gamma_{n+1}(t)\geqslant \gamma_n(t)+\ln|t|_p-\ln p$ $\forall


t\in\Bbb C^p$. Следовательно, топология в $C^\infty(F)$\; $d$-разделена с


$d(t)=\ln|t|_p$.


Предположим, что $\Lambda=(t_{(k)})$ имеет конечный логарифмический


$d$-тип, т.\,е.


\begin{equation}


\limsup\limits_{k\to\infty}\dfrac{\ln(1+k)}{\ln|t_{(k)}|_p}


<\infty.


\end{equation}


Полагая $a_k=\ln|t_{(k)}|_p$, по теореме 4 заключаем, что ряд (21) сходится


(или абсолютно сходится) в $C^\infty(F)$ тогда и только тогда, когда


$$


\limsup\limits_{k\to\infty}\dfrac{1}{\ln|t_{(k)}|_p}\bigl(


\ln|c_k|+\gamma_n(t_{(k)})+h_F(t_{(k)})\bigr)\leqslant 0


\ \; \forall n\geqslant 1.


$$


При этом


$$


\gamma_n(t_{(k)})\geqslant\gamma_{0}(t_{(k)})+


n\ln|t_{(k)}|_p-n\ln p=1+n\ln|t_{(k)}|_p-n\ln p


\ \; \forall k\geqslant 1.


$$


Следовательно, если ряд (21) сходится в $C^\infty(F)$, то


$$


\limsup\limits_{k\to\infty}\dfrac{1}{\ln|t_{(k)}|_p}\bigl(


\ln|c_k|+h_F(t_{(k)})\bigr)\leqslant -n


\ \; \forall n\geqslant 1,


$$


т.\,е.


\begin{equation}


\lim\limits_{k\to\infty}\dfrac{1}{\ln|t_{(k)}|_p}\bigl(


\ln|c_k|+h_F(t_{(k)})\bigr)=-\infty.


\end{equation}


С другой стороны, $\gamma_n(t)\leqslant


n\ln|t_{(k)}|_p$ $\forall n\geqslant 1$.


Поэтому, если равенство (32) имеет место, то


$$


\lim\limits_{k\to\infty}\dfrac{1}{\ln|t_{(k)}|_p}\bigl(


\ln|c_k|+\gamma_n(t_{(k)})+h_F(t_{(k)})\bigr)=-\infty


\ \; \forall n\geqslant 1.


$$


Таким образом, ряд (21) при условии (31) сходится


(или абсолютно сходится) в $C^\infty(F)$ тогда и только тогда, когда


выполнено условие (32).





Автор выражает благодарность рецензенту за полезные замечания.





\begin{thebibliography}{9}





\bibitem{1}


Коробейник Ю.Ф. {\em Об одной двойственной задаче}. I. {\em Общие результаты.


Приложения к пространствам Фреше} // Матем. сб. -- 1975. --  Т.\,97. --


\No 2. -- С.\,193--229.


\bibitem{2}


Маркушевич А.И. {\em Теория аналитических функций}. -- М.: ГИТТЛ, 1950. --


703 c.





\end{thebibliography}





\vskip 2cm





\post{\it Ростовский государственный университет}{\it Поступила}


\post{}{28.06.1999}





\label{end}


\end{document}


